% !TeX program = xelatex
\documentclass[11pt,a4paper]{article}
\usepackage{fontspec,polyglossia,amsmath,amssymb,booktabs,geometry}
\setdefaultlanguage{ukrainian}
\setmainfont{DejaVu Serif}
\newfontfamily\cyrillicfont{DejaVu Serif}
\geometry{margin=22mm}
\setlength\parindent{0pt}
\setlength{\emergencystretch}{3em}
\setlength\parskip{4pt}
\raggedbottom
\begin{document}
{\small Обчислювальна математика та програмування \hfill Лекція 7}\\[4mm]
{\LARGE Диференціювання та інтегрування функцій}\\[2mm]
Математичний матеріал до лабораторної роботи №6
\section{Швидкість і прискорення за формулою руху}
Для внутрішньої точки $x_0$ області визначення похідна існує, якщо існує скінченна границя
\[
 f'(x_0)=\lim_{h\to0}\frac{f(x_0+h)-f(x_0)}{h}.
\]
Графік дотичної має рівняння $y=f(x_0)+f'(x_0)(x-x_0)$.
Друга похідна $f''(x_0)=(f')'(x_0)$ існує за диференційовності $f'$ у цій точці.
Для координати $s(t)$, заданої в метрах, і часу $t$ у секундах: $v=s'$ (м/с), $a=v'=s''$ (м/с$^2$).
У моделі $s(t)=t^2$ з числовими значеннями в цих одиницях
\[
\frac{s(t+h)-s(t)}h=2t+h,\quad v(t)=2t,\quad a(t)=2.
\]
\section{Похідні за сусідніми вимірами}
Нехай $h>0$, $f_j=f(x_0+jh)$, а всі використані вузли належать області визначення. У таблицях нижче нумерація відповідає табл.~6.1 та 6.2 методички. Для оцінок порядку похибки припускається достатня гладкість $f$ на відрізку використаних вузлів.
\begin{center}\begin{tabular}{c l l}
\toprule № & Наближення першої похідної & Похибка апроксимації\\\midrule
1 & $(f_1-f_0)/h$ & $O(h)$\\
2 & $(-3f_0+4f_1-f_2)/(2h)$ & $O(h^2)$\\
3 & $(f_0-f_{-1})/h$ & $O(h)$\\
4 & $(3f_0-4f_{-1}+f_{-2})/(2h)$ & $O(h^2)$\\
5 & $(f_1-f_{-1})/(2h)$ & $O(h^2)$\\
6 & $(f_{-2}-8f_{-1}+8f_1-f_2)/(12h)$ & $O(h^4)$\\\bottomrule
\end{tabular}\end{center}
\begin{center}\begin{tabular}{c l l}
\toprule № & Наближення другої похідної & Похибка апроксимації\\\midrule
1 & $(f_2-2f_1+f_0)/h^2$ & $O(h)$\\
2 & $(2f_0-5f_1+4f_2-f_3)/h^2$ & $O(h^2)$\\
3 & $(f_0-2f_{-1}+f_{-2})/h^2$ & $O(h)$\\
4 & $(2f_0-5f_{-1}+4f_{-2}-f_{-3})/h^2$ & $O(h^2)$\\
5 & $(f_1-2f_0+f_{-1})/h^2$ & $O(h^2)$\\
6 & $(-f_2+16f_1-30f_0+16f_{-1}-f_{-2})/(12h^2)$ & $O(h^4)$\\\bottomrule
\end{tabular}\end{center}
За еталона $D$ та наближення $\widehat D$ абсолютна похибка $e=|\widehat D-D|$.
Відносна похибка $e/|D|$ визначена при $D\ne0$.
Якщо $\widehat D=h^{-m}\sum_j w_j f_j$ і $|\delta f_j|\le\varepsilon$, то
\[
 |\delta\widehat D|\le h^{-m}\varepsilon\sum_j|w_j|.
\]
Для центральної другої похідної із єдиним зсувом $\delta f_0=\eta$ маємо
$\delta\widehat D=-2\eta/h^2$. Для $s(t)=t^3$ та $t_0=2$ права й центральна перші різниці дорівнюють $12+6h+h^2$ та $12+h^2$ відповідно; центральна друга різниця дорівнює $12$.
\section{Накопичення і визначений інтеграл}
Функція $F$ є первісною $q$ на інтервалі, якщо $F'=q$. Для неперервної $q$ на $[a,b]$:
\[
 \int q(t)\,dt=F(t)+C,\qquad \int_a^b q(t)\,dt=F(b)-F(a).
\]
При $a<b$, $t_i=a+i\Delta t$, $\Delta t=(b-a)/n$, $\xi_i\in[t_i,t_{i+1}]$ маємо
\[
 \int_a^b q(t)\,dt=\lim_{n\to\infty}\sum_{i=0}^{n-1}q(\xi_i)\Delta t .
\]
Вибір $\xi_i=t_i$, $t_{i+1}$ або $(t_i+t_{i+1})/2$ задає ліві, праві або середні прямокутники. За $q\in C^2[a,b]$ похибка складених середніх прямокутників є $O(\Delta t^2)$, лівих або правих у загальному випадку $O(\Delta t)$.
Нехай $q(t)=2t-4$, $0\le b\le4$, $V(0)=10$. Тоді
\[
 \Delta V=b^2-4b,\quad V(b)=10+b^2-4b,\quad
 \int_0^b |q(t)|\,dt=
 \begin{cases}4b-b^2,&b\le2,\\4+(b-2)^2,&b>2.\end{cases}
\]
\section{Числове інтегрування та перевірка похибки}
Для $n$ рівних інтервалів, $h=(b-a)/n$, $P_i=P(a+ih)$:
\[
 T_n=h\left(\frac{P_0+P_n}{2}+\sum_{i=1}^{n-1}P_i\right).
\]
Якщо $P\in C^2[a,b]$, то
$\left|T_n-\int_a^b P\right|\le (b-a)h^2\max_{[a,b]}|P''|/12$.
Для модельної функції з $A>0$, $w>0$ і центром $c$:
\[
 P(t)=\frac{A}{1+((t-c)/w)^2},\qquad
 I=Aw\left[\arctan\frac{b-c}{w}-\arctan\frac{a-c}{w}\right].
\]
Для оцінки квадратури $Q$ її фактична абсолютна похибка становить $|Q-I|$. Внутрішня оцінка $\widehat e$ залежить від алгоритму й загалом не збігається з $|Q-I|$.
Абсолютний і відносний допуски $\varepsilon_a,\varepsilon_r$ задають цільову точність $\max(\varepsilon_a,\varepsilon_r|I|)$. Для $P$ у ватах і $t$ у годинах $I$, $Q$, $|Q-I|$ та $\widehat e$ мають розмірність Вт$\cdot$год.
\end{document}
